\section{Extracting a prescribed cusp count}\label{extracting-fixed-cusp-counts}

The preceding estimates control generating functions, which sum over cusp counts. We now recover a single coefficient, and hence a volume with a prescribed number of cusps. A positive convolution factor supplies Fourier bounds that turn the real-axis estimates into a local limit theorem. Small cusp weight, the interior, and the critical endpoint require separate estimates, with uniform constants across the three ranges.

Subsection~\ref{canonical-coefficient-ratios-and-all-cusp-count-uniformity} proves the qualitative volume-ratio limits using the leading estimates from Sections~\ref{ratios-with-a-fixed-number-of-distinguished-markings}--\ref{sec:prefactor}. The final subsection extends the argument to bounded geodesic boundaries. The higher-order expansions are needed only for the sharper extraction estimates that follow.

\subsection{A positive convolution factor}\label{a-positive-convolution-factor}

\begin{lemma}[Positive convolution factor]\label{lem:convolution}
For each stable marked pair \((g,p)\) with \(p\ge1\), the marked generating
function has the exact factorization
\begin{equation}
\Phi_g^{(p)}(x)=V(x)^h R_{g,p}(x),\qquad
 V(x)=Y'(x)=1/a(Y(x)),\qquad h=2g-2+p,
 \label{eq:TL33}
\end{equation} where \(R_{g,p}\) has nonnegative Taylor coefficients and
converges for \(0\le x<x_c\).
\end{lemma}

The point of \eqref{eq:TL33} is the following. An approximation of \(\Phi_g^{(p)}\) on
the real axis alone says nothing about its individual coefficients. The
factorization shows that, under the cusp weight \(x\), the number of
additional cusps is distributed as the sum of \(h\) independent variables
with generating function \(V(xz)/V(x)\) and one further independent
nonnegative integer variable; this gives the Fourier bound \eqref{eq:TL35} below, which is what
a local limit theorem needs.

\begin{proof}
The proof uses positivity in the parameter: by
Lemma~\ref{lem:kernel-positivity}, the functions
\(\widehat\vartheta_\ell(Y)=a(Y)[z^\ell]K_Y(z)^{-1}\) and \(1/a(Y)\) have
nonnegative Taylor coefficients in \(Y\). Since this positivity in \(Y\) is
available for the unrescaled coefficients, we do not use the rescaling by
\(r(Y)\) of \eqref{eq:H7} here.

Undo the normalization in \eqref{eq:H9} and regard the unrescaled coefficients \(\vartheta_\ell=[z^\ell]K_Y(z)^{-1}\) as independent variables. In this convention the bases are \(U_{0,3}=\vartheta_0\), \(U_{1,1}(1)=\vartheta_0/8\), and \(U_{1,1}(0)=\vartheta_1/8\). Induction from these bases shows that every amplitude is a finite polynomial with nonnegative
coefficients in the variables \(\vartheta_\ell\); each monomial has total degree \(h\),
and weighted degree \(m\) when \(\vartheta_\ell\) is given weight \(\ell\). A
linear recursion term adds one factor \(\vartheta_\ell\) to a child of complexity \(h-1\); a
separating term adds one to children whose complexities sum to \(h-1\).
Thus factoring \(a^{-1}\) from each of the \(h\) factors \(\vartheta_\ell=a^{-1}\widehat\vartheta_\ell\) leaves a
positive polynomial in the \(\widehat\vartheta_\ell\).

Finally \(Y(x)\) has nonnegative coefficients and \(Y(0)=0\). To justify
the product used here, a nonzero zero \(y\) of \(X\) gives
\(u(s)=J_1(2\sqrt y\,s)\), with \(u(1)=0\) and regular behavior at zero.
Integration of \((s u')'+(4ys-1/s)u=0\) against its conjugate gives
\(4y\int_0^1s|u|^2\,\mathrm{d}s=\int_0^1(s|u'|^2+|u|^2/s)\,\mathrm{d}s>0\). Thus all
nonzero zeros are positive real. Since \(X\) has order one half and a
simple zero at zero with derivative one, its Hadamard product is
\(X(Y)=Y\prod_k(1-Y/\beta_k)\), \(\beta_k>0\). Lagrange inversion now
gives \[
 [x^k]Y(x)=\frac1k[Y^{k-1}](Y/X(Y))^k>0.
\] The inverse series has radius exactly \(x_c\): a smaller radius
would, by positivity and Pringsheim's theorem, give a positive real
singularity before \(x_c\), contradicting the analytic real inverse
there; the quadratic critical point at \(x_c\) supplies the opposite
inequality. Composition therefore preserves coefficient positivity and
converges for \(0\le x<x_c\), since \(Y(x)<y_0\). Finally, by
\eqref{eq:H8}, \(\Phi_g^{(p)}(x)=U^{Y}_{g,p}(0^p)\) with \(Y=Y(x)\), and the
factor \(a(Y)^{-h}\) removed above equals \(V(x)^h\). This proves
\eqref{eq:TL33}.
\end{proof}

For the coefficient probability law with generating function
\(\Phi_g^{(p)}(xz)/\Phi_g^{(p)}(x)\), positivity in \eqref{eq:TL33} gives \begin{equation}
\left|\frac{\Phi_g^{(p)}(x\mathrm{e}^{i\theta})}{\Phi_g^{(p)}(x)}\right|
 \le\left|\frac{V(x\mathrm{e}^{i\theta})}{V(x)}\right|^h.
 \label{eq:TL35}
\end{equation} This is the Fourier estimate that prevents a generating-function
ratio from concealing different local coefficient behavior.

\subsection{Mean, variance and a central limit theorem}\label{mean-variance-and-a-real-central-limit-theorem}

We first compute the mean and variance of the cusp count under the cusp
weight \(x\), and show that it is asymptotically Gaussian; the next
subsection upgrades this to a local limit theorem. Recall
\(\mathscr D=x\,\mathrm{d}/\mathrm{d}x\) and the function \(\Psi\) of
Lemma~\ref{lem:parameter}. For convenience we repeat its endpoint expansions
\eqref{eq:SP16}, together with that of \(V=1/a(Y(x))\), which follows from
\(a(Y)=1-Y+O(Y^2)\) and \(Y(x)=x+O(x^2)\): \begin{equation}
\begin{aligned}
 x\downarrow0:\quad& V(x)=1+x+O(x^2),\quad \Psi(x)=2x+O(x^2),\\
 &\mathscr D\Psi(x)=2x+O(x^2),\quad \mathscr D^2\Psi(x)=2x+O(x^2);\\[1mm]
 x\uparrow x_c:\quad&\Psi(x)=\frac5{4\varepsilon}(1+O(\sqrt\varepsilon)),\\
 &\mathscr D\Psi(x)=\frac5{4\varepsilon^2}(1+O(\sqrt\varepsilon)),\quad
 \mathscr D^2\Psi(x)=O(\varepsilon^{-3}),
 \qquad\varepsilon=1-x/x_c.
 \end{aligned}
 \label{eq:TL36}
\end{equation} On compact subintervals of \((0,x_c)\), all needed
derivatives are bounded and \(\mathscr D\Psi\) is bounded away from
zero.

For \(F=\Phi_g^{(p)}\) at a positive cusp weight \(x\), let \begin{equation}
\mathbb P_x(N=k)=\frac{[z^k]F(z)x^k}{F(x)},\qquad
 \mu=h\Psi(x),\qquad \sigma^2=h\mathscr D\Psi(x).
 \label{eq:TL37}
\end{equation} The last two quantities are model center and variance, not claims
about their exact finite-genus values. Suppose \(h\to\infty\) and
\(\mu\to\infty\). For every fixed real \(u\), integrating \eqref{eq:TL30} gives
\begin{equation}
\log\mathbb E_x \mathrm{e}^{u(N-\mu)/\sigma}
 =\int_0^{u/\sigma}
 \{h\Psi(x\mathrm{e}^s)-h\Psi(x)+O_p(\Psi(x\mathrm{e}^s))\}\,\mathrm{d}s.
 \label{eq:TL38}
\end{equation} The interval of integration stays inside \((0,x_c)\). Expanding
\(h\Psi(x\mathrm e^s)-h\Psi(x)=sh\mathscr D\Psi(x)+O(s^2h\mathscr D^2\Psi)\),
the first part of the integral is \(u^2/2\) plus a Taylor remainder, and the
second part is an integrated error; both are \(o(1)\).
More precisely, in the critical regime its Taylor remainder and
integrated error are \(O_u(h^{-1/2})\); on an interior compact interval
they are also \(O_u(h^{-1/2})\). At \(x\to0\), \(\sigma^2\sim\mu\), the
Taylor remainder is \(O_u(\mu^{-1/2})\) and the integrated error is
\(O_u(x/\sqrt\mu)=O_u(\sqrt\mu/h)=o(1)\). These assertions follow
directly from \eqref{eq:TL36}; no estimate for the derivative of the
error in \eqref{eq:TL30} is required.

Thus the real moment generating functions converge to \(\exp(u^2/2)\)
along every such sequence. Exponential Markov bounds give tightness;
bounds at slightly larger real arguments give uniform integrability
along weakly convergent subsequences. Their limiting moment generating
functions identify the standard Gaussian. Therefore \begin{equation}
(N-\mu)/\sigma\Longrightarrow\mathcal N(0,1).
 \label{eq:TL39}
\end{equation}

\subsection{Fourier bounds and the local limit theorem}\label{uniform-fourier-bounds-and-the-lattice-local-limit}

A central limit theorem alone does not determine ratios of individual
coefficients. We now prove the corresponding local limit theorem.

\begin{proposition}[Uniform local limit theorem]\label{prop:local-limit}
Fix \(p\ge1\) and use the notation \eqref{eq:TL37}. Along every sequence
with \(h\to\infty\), \(\mu\to\infty\) and \(x\in(0,x_c)\),
\begin{equation}
\mathbb P_x(N=k)=\frac{1+o(1)}{\sqrt{2\pi}\sigma}
 \mathrm{e}^{-u^2/2}\qquad\text{whenever }(k-\mu)/\sigma\longrightarrow u.
 \label{eq:TL47}
\end{equation}
No relation between \(h\) and the distance from \(x\) to \(x_c\) is required.
\end{proposition}

\begin{proof}
Passing to a subsequence, we may assume that \(x\) tends to zero, to an
interior point, or to \(x_c\). In each case we bound the characteristic
function of \((N-\mu)/\sigma\) by an integrable function, uniformly along
the sequence; Fourier inversion and the pointwise convergence implied by
\eqref{eq:TL39} then give \eqref{eq:TL47}.

On a compact subinterval of \((0,x_c)\), the coefficient probabilities
at zero and one of \(V(xz)/V(x)\) have positive uniform lower bounds. If
these probabilities are \(p_0,p_1\), then \begin{equation}
\left|\sum_jp_j\mathrm{e}^{ij\theta}\right|^2
 \le1-2p_0p_1(1-\cos\theta).
 \label{eq:TL40}
\end{equation} Equations \eqref{eq:TL35} and \eqref{eq:TL40} give an upper bound
\(\exp(-c h\theta^2)\) on \([-\pi,\pi]\). Since \(\sigma\) is comparable
to \(\sqrt h\), this is integrable Gaussian domination on the rescaled
Fourier axis.

At \(x\to0\), one has \(p_0\asymp1\), \(p_1\asymp x\). The same argument
gives \(\exp(-c h x\theta^2)\). Since \(hx\asymp\mu\) and
\(\sigma^2\sim\mu\to\infty\), it again gives integrable domination after
rescaling.

Near the critical point a different bound is needed, because the law
with generating function \(V(xz)/V(x)\) then has variance of order
\(\varepsilon^{-2}\) and its probabilities at zero and one tend to zero. Quadratic
analytic inversion at \(x_c\) gives \begin{equation}
V(z)=A_0(1-z/x_c)^{-1/2}
 \Xi\bigl(\sqrt{1-z/x_c}\bigr),\qquad
 A_0>0,\quad \Xi(0)=1,
 \label{eq:TL41}
\end{equation} where \(\Xi\) is analytic and nonzero near zero and has real
coefficients. Put \(\Lambda(w)=\log \Xi(\sqrt w)\) and
\(w(\theta)=1-(1-\varepsilon)\mathrm{e}^{i\theta}\). Then \begin{equation}
\log\left|\frac{V(x\mathrm{e}^{i\theta})}{V(x)}\right|
 =-\frac14\log\left(1+
 \frac{2(1-\varepsilon)(1-\cos\theta)}{\varepsilon^2}\right)
 +\Re \Lambda(w(\theta))-\Lambda(\varepsilon).
 \label{eq:TL42}
\end{equation} There are fixed \(\delta,\varepsilon_0>0\) such that \begin{equation}
\left|\frac{V(x\mathrm{e}^{i\theta})}{V(x)}\right|
 \le(1+\theta^2/(4\varepsilon^2))^{-1/8}
 \quad(|\theta|\le\delta,\ 0<\varepsilon\le\varepsilon_0).
 \label{eq:TL43}
\end{equation} Indeed, for \(|\theta|\le\varepsilon\), the first derivative at zero
of the real part of \(\Lambda(w(\theta))\) vanishes, while its second
derivative is \(O(\varepsilon^{-3/2})\). Its change is therefore
\(O(\theta^2\varepsilon^{-3/2})\), absorbed by half of the leading
logarithm in \eqref{eq:TL42} when \(\varepsilon\) is small. For
\(\varepsilon\le|\theta|\le\delta\), its change is bounded by
\(C\sqrt\delta\), while the logarithm is at least \(\log(5/4)\). Choose
\(\delta\) small enough for the same absorption. We used the elementary
bound \(2(1-\varepsilon)(1-\cos\theta)\ge\theta^2/4\) on the small arc.

On the outer arc \(\delta\le|\theta|\le\pi\), \begin{equation}
|V(x\mathrm{e}^{i\theta})/V(x)|\le C_\delta\sqrt\varepsilon.
 \label{eq:TL44}
\end{equation} To see that \(V\) is analytic there, note that nonnegative coefficients and
\(Y(x)\to y_0\) give absolute convergence of \(Y\) on \(|x|\le x_c\).
Its first two coefficients are positive, so strict triangle inequality
gives \(|Y(x_c \mathrm{e}^{i\theta})|<y_0\) away from \(\theta=0\). By Lemma~\ref{lem:kernel-zeros}, \(X'\) has no zero in \(|Y|<y_0\). The inverse-function theorem
extends \(Y,V\) across the outer arc, where \(V\) is uniformly bounded.
Since \eqref{eq:TL41} gives \(V(x)\asymp\varepsilon^{-1/2}\), \eqref{eq:TL44} follows.

On the critical Fourier scale \(\theta=\varepsilon t/\sqrt h\), \eqref{eq:TL35}
and \eqref{eq:TL43} are bounded by \begin{equation}
(1+t^2/(4h))^{-h/8}\le(1+t^2/64)^{-2},\qquad h\ge16.
 \label{eq:TL45}
\end{equation} The inequality follows because \((1+a/h)^h\) increases with \(h\).
The right side is integrable. The outer contribution after
multiplication by \(\sqrt h/\varepsilon\) is at most \begin{equation}
2\pi\sqrt h\,C_\delta^h\varepsilon^{h/2-1}\longrightarrow0.
 \label{eq:TL46}
\end{equation} This convergence is uniform for arbitrarily small \(\varepsilon\):
fix sufficiently small \(\varepsilon_0\) with
\(C_\delta\sqrt{\varepsilon_0}<1\), and bound \eqref{eq:TL46} by
\(2\pi\varepsilon_0^{-1}\sqrt h(C_\delta\sqrt{\varepsilon_0})^h\) for
\(h\ge2\).

These integrable bounds, in the three regimes, complete the proof.
\end{proof}

The same Fourier bounds give an upper bound for every point mass, not
only near the center. We record it for later use in
Section~\ref{sec:geometry} and Appendix~\ref{app:tight}.

\begin{lemma}[Uniform point-mass bound]\label{lem:point-mass}
For every $c_0>0$ there is a constant $C$ with the following property.
Let $x\in(0,x_c)$, let $h,h'$ be positive integers with $h'\ge\max(16,c_0h)$,
and let $G(z)=V(z)^{h'}R(z)$, where $R$ has nonnegative Taylor
coefficients, converges for $|z|<x_c$ and satisfies $R(x)>0$. Then,
with $\sigma^2=h\mathscr D\Psi(x)$,
\[
 \sup_{j\ge0}\frac{[z^j]G(z)\,x^j}{G(x)}\le\frac C\sigma.
\]
\end{lemma}

\begin{proof}
By Fourier inversion and positivity of the coefficients of $R$,
\[
 \frac{[z^j]G(z)\,x^j}{G(x)}
 \le\frac1{2\pi}\int_{-\pi}^{\pi}
       \left|\frac{V(x\mathrm e^{i\theta})}{V(x)}\right|^{h'}\mathrm d\theta.
\]
On a compact subinterval of $(0,x_c)$, \eqref{eq:TL40} bounds the
integrand by $\mathrm e^{-ch'\theta^2}$, so the integral is
$O(h^{-1/2})$, while $\sigma\asymp\sqrt h$. Near zero the same argument
with $p_1\asymp x$ gives $\mathrm e^{-ch'x\theta^2}$ and the bound
$O((hx)^{-1/2})$, while $\sigma^2\asymp hx$. Near $x_c$, put
$\theta=\varepsilon t/\sqrt{h'}$ on $|\theta|\le\delta$; \eqref{eq:TL43}
and the argument of \eqref{eq:TL45} bound this part by
$(\varepsilon/\sqrt{h'})\int(1+t^2/64)^{-2}\,\mathrm dt$. On the outer arc
\eqref{eq:TL44} gives at most
$2\pi(C_\delta\sqrt\varepsilon)^{h'}\le C\varepsilon h^{-1/2}$ once
$C_\delta\sqrt{\varepsilon_0}<1$. Since $\sigma\asymp\sqrt h/\varepsilon$
there, all three regimes give $O(1/\sigma)$.
\end{proof}

\subsection{Ratios at a prescribed number of cusps}\label{canonical-coefficient-ratios-and-all-cusp-count-uniformity}

We now prove the limits in Theorem~\ref{thm:volume}. Each volume ratio is
a ratio of coefficients, at the same index, of two marked generating
functions. We evaluate both at the same cusp weight, chosen so that the
index is, up to a negligible shift, the model center of both. The local limit theorem \eqref{eq:TL47} shows
that the two local probabilities are asymptotically equal, so the
coefficient ratio is asymptotic to the ratio of the generating functions
themselves, which Section~\ref{ratios-with-a-fixed-number-of-distinguished-markings}
determines.

\begin{proof}[Uniform convergence in Theorem~\ref{thm:volume}]
First suppose \(g,n\to\infty\). Put \begin{equation}
k=n-1,\qquad h=2g-1,
 \qquad h\Psi(x)=k.
 \label{eq:TL48}
\end{equation} The last equation has a unique solution in \((0,x_c)\) by \eqref{eq:TL32} and
the parameter bijection. The two families \(\Phi_g'\) and
\(\Phi_{g-1}'''\) have exactly the same complexity \(h\), model center,
and model variance. Thus their two local probabilities in \eqref{eq:TL47}, at the
common coefficient index \(k\), have ratio tending to one. Equations
\eqref{eq:TL31} and \eqref{eq:TL47} give \begin{equation}
\frac{[x^k]\Phi_{g-1}'''(x)}{[x^k]\Phi_g'(x)}
 =4r(Y)(1+o(1)).
 \label{eq:TL49}
\end{equation} By the definition \eqref{eq:intro-Phi}, the coefficients are \begin{equation}
[x^k]\Phi_g'(x)=
 \frac{V_{g,n}}{(2\pi^2)^{3g-3+n}k!},\qquad
 [x^k]\Phi_{g-1}'''(x)=
 \frac{V_{g-1,n+2}}{(2\pi^2)^{3g-4+n}k!}.
 \label{eq:TL50}
\end{equation} Consequently \begin{equation}
B_{g,n}=\frac{\pi^2}{2r(Y)}(1+o(1))
 =\mathcal B\left(\frac{2(n-1)}{2g-1}\right)(1+o(1)).
 \label{eq:TL51}
\end{equation} The factorial \(k!\) cancels, and the powers of \(2\pi^2\) differ by
exactly one; this explains the factor in \eqref{eq:TL51}.

For the cusp ratio, compare \(\Phi_g''\) with \(\Phi_g'\) at the same
index \(k\) and the same cusp weight \eqref{eq:TL48}. Their complexities are \(h+1\)
and \(h\). Their centers differ by \(\Psi(x)\), and their variances have
ratio \((h+1)/h\). Moreover \begin{equation}
\frac{\Psi(x)}{\sqrt{h\mathscr D\Psi(x)}}=O(h^{-1/2})
 \label{eq:TL52}
\end{equation} uniformly in \(x\), because \(\Psi^2/(\mathscr D\Psi)\) is continuous
in the interior and has finite endpoint limits zero and \(5/4\),
respectively. Thus the local probability ratio again tends to one.
Undoing \eqref{eq:TL24}, \begin{equation}
C_{g,n}=\frac{\pi^2}{g}
 \frac{[x^k]\Phi_g''}{[x^k]\Phi_g'}
 =\frac{\pi^2}{g}\frac{h}{Aar}(1+o(1))
 =\frac{\pi^2(n-1)}{gx}(1+o(1)).
 \label{eq:TL53}
\end{equation} Here, as in \eqref{eq:TL50}, the coefficients \([x^k]\Phi_g''\) and
\([x^k]\Phi_g'\) carry the normalizing powers \((2\pi^2)^{3g-2+n}\) and
\((2\pi^2)^{3g-3+n}\); division by \(2g\) gives the factor \(\pi^2/g\). With
\(\lambda_{\mathrm{eff}}=2(n-1)/(2g-1)\), the last expression equals \begin{equation}
\left(1-\frac1{2g}\right)\mathcal C(\lambda_{\rm eff})(1+o(1)).
 \label{eq:TL54}
\end{equation}

One may replace \(\lambda_{\mathrm{eff}}\) by \(n/g\) in both \eqref{eq:TL51} and
\eqref{eq:TL54}. At a zero limiting slope, use the continuous values
\(\mathcal B(0)=1\) and \(\mathcal C(0)=4\pi^2\) of \eqref{eq:SP3}; at a
finite positive slope use continuity; and at an infinite slope use the
endpoint equivalents \eqref{eq:SP15} of Lemma~\ref{lem:parameter}.
Since \(\lambda_{\mathrm{eff}}/(n/g)\to1\) whenever \(n\to\infty\),
these give the desired relative replacement in all cases.

Finally the exact identity \begin{equation}
D_{g,n}=C_{g,n-1}/C_{g,n}
 \label{eq:TL56}
\end{equation} and the adjacent continuity of the profile \(\mathcal C\) give
\(D_{g,n}\to1\) for every growing-cusp sequence. For bounded cusp
counts, the classical fixed-cusp theorems give \(B_{g,n}\to1\),
\(C_{g,n}\to4\pi^2\), and hence \(D_{g,n}\to1\); these are the only
bounded-cusp inputs used here.

If any supremum in Theorem~\ref{thm:volume} did not tend to zero, select a violating
sequence of cusp counts. A subsequence either has a fixed bounded cusp
count or tends to infinity. The preceding conclusions contradict a fixed
violation in both cases. This proves the uniform convergence assertions.
Section~\ref{subsec:quantitative-volume-ratios} supplies the stated
inverse-genus rate using the sharper extraction theorem.

\end{proof}

\subsection{Bounded geodesic boundaries}\label{subsec:bounded-boundaries}

The positive factor also permits a fixed number of cusps to be replaced
by geodesic boundaries. The resulting factor depends on the cusp density
and joins the familiar large-genus and large-cusp boundary profiles.
For $0\le Y\le y_0$ and $L\ge0$, put
\begin{equation}
 \mathfrak b_Y(L)=\int_0^1
 I_0\!\left(\frac L\pi\sqrt{T(Y)-\Delta(Y)s^2}\right)\,\mathrm ds,
 \label{eq:boundary-profile}
\end{equation}
where $I_0$ is the modified Bessel function of order zero, $T=Y+\Delta$
and $r=2\Delta$. At $Y=y_0$ use $T=y_0$ and $\Delta=0$.
Writing $Y_\lambda$ for the inverse of $\lambda(Y)$, we have
$Y_0=0$ and $Y_\infty=y_0$.

\begin{proposition}[Boundary lengths at arbitrary cusp density]
\label{prop:bounded-boundaries}
Fix an integer $p\ge1$ and $L_0<\infty$. As $g\to\infty$,
\begin{equation}
 \sup_{\substack{n\ge p\\\mathbf L\in[0,L_0]^p}}
 \left|
 \frac{V_{g,n}(L_1,\ldots,L_p,0^{n-p})}
 {V_{g,n}\prod_{i=1}^p\mathfrak b_{Y_{n/g}}(L_i)}-1
 \right|\longrightarrow0.
 \label{eq:boundary-uniform}
\end{equation}
The endpoint factors are
\begin{equation}
 \mathfrak b_0(L)=\frac{\sinh(L/2)}{L/2},\qquad
 \mathfrak b_{y_0}(L)=I_0\!\left(\frac{j_{0,1}L}{2\pi}\right),
 \label{eq:boundary-endpoints}
\end{equation}
with value one at $L=0$.
\end{proposition}

The first factor occurs in the fixed-cusp estimate of Mirzakhani and
Petri~\cite[Proposition~3.1]{MP}; the second is the fixed-genus,
large-cusp factor of Hide and Thomas~\cite[Theorem~1.5]{HTLargeN}.
Here the genus tends to infinity and the cusp count is unrestricted.
The number and lengths of the nonzero boundaries remain bounded.
For comparison, Budd and Lions obtain a boundary factor for tight volumes
with a random cusp count in a near-critical regime
\cite[Proposition~4.2]{BL}. Exact transformations between ordinary and
tight boundary generating functions were established in~\cite{BZ}.

\begin{proof}
Put $h=2g-2+p$, $F(x)=\Phi_g^{(p)}(x)$ and
\[
 G(x;\mathbf L)=\sum_{k\ge0}
 \frac{V_{g,p+k}(\mathbf L,0^k)}{(2\pi^2)^{3g-3+p+k}}
 \frac{x^k}{k!}.
\]
We first compare $G$ with $F$ on the real axis, and then show that the
same comparison holds for their coefficients at the saddle.

\emph{The boundary insertion.}
Use the string flow of \eqref{eq:H8}, holding its parameter $Y$ fixed
while differentiating the KW times. Its linear part gives
\[
 \partial_{t_d}=\sum_{b=0}^d\frac{Y^{d-b}}{(d-b)!}\partial_{t_b'}.
\]
For $g\ge2$ the free energy is invariant under the flow. Apply the
displayed identity to each external insertion before setting $x=X(Y)$.
In the volume normalization of \eqref{eq:H5}, an external degree $d$
has coefficient $(L^2/(4\pi^2))^d/d!$ times the unweighted derivative.
Consequently, with $t_i=L_i^2/(4\pi^2)$, the exact identity is
\begin{equation}
 \frac{G(X(Y);\mathbf L)}{F(X(Y))}
 =\sum_{\mathbf a,\mathbf b\ge0}c_{g,p}(\mathbf b;Y)
 \prod_{i=1}^p
 \frac{t_i^{a_i+b_i}Y^{a_i}r^{b_i}}
 {a_i!(a_i+b_i)!(2b_i+1)!!}.
 \label{eq:boundary-insertion}
\end{equation}
Indeed, \eqref{eq:TL9} gives
$U^Y_{g,p}(\mathbf b)/U^Y_{g,p}(\mathbf0)=r^{|\mathbf b|}c_{g,p}(\mathbf b;Y)$.
Terms with invalid $\mathbf b$ are zero; there is no bound on $\mathbf a$.

The identity
\[
 \int_0^1(1-s^2)^b\,\mathrm ds=\frac{2^b b!}{(2b+1)!!}
\]
and the series for $I_0$ show that replacing $c_{g,p}$ by one in
\eqref{eq:boundary-insertion} gives
$\prod_i\mathfrak b_Y(L_i)$. The same double series with an extra
factor $|\mathbf b|(|\mathbf b|+1)$ is bounded uniformly for
$Y\in[0,y_0]$ and $\mathbf L\in[0,L_0]^p$: the inequality
$(a+b)!\ge a!b!$ supplies a summable factorial majorant for every fixed
polynomial moment. Thus \eqref{eq:TL23}, including its one-marking case,
gives
\begin{equation}
 \frac{G(X(Y);\mathbf L)}{F(X(Y))}
 =\prod_i\mathfrak b_Y(L_i)\bigl(1+O_{p,L_0}(g^{-1})\bigr).
 \label{eq:boundary-real}
\end{equation}
All terms in \eqref{eq:boundary-insertion} are nonnegative. The
majorant also justifies convergence of the formal string-flow identity
for $0\le x<x_c$. The factors $\mathfrak b_Y(L)$ are at least one;
they and their first $Y$ derivatives are uniformly bounded on the
compact set under consideration.

\emph{Extraction at a prescribed cusp count.}
The proof of \eqref{eq:TL33} applies to every external degree tuple:
$a(Y)^hU^Y_{g,p}(\mathbf b)$ has nonnegative Taylor coefficients in $Y$.
Before division by $F$, the insertion identity is a sum of these
amplitudes with nonnegative powers of $Y$ and nonnegative constants.
Since $Y(x)$ also has nonnegative coefficients, it follows that
\[
 G(x;\mathbf L)=V(x)^h R_{g,p,\mathbf L}(x),\qquad
 [x^j]R_{g,p,\mathbf L}(x)\ge0.
\]
Convergence follows from \eqref{eq:boundary-real}. Hence
\begin{equation}
 \left|\frac{G(x\mathrm e^{i\theta};\mathbf L)}{G(x;\mathbf L)}\right|
 \le\left|\frac{V(x\mathrm e^{i\theta})}{V(x)}\right|^h.
 \label{eq:boundary-fourier}
\end{equation}

Suppose first that $k=n-p\to\infty$, and choose $x$ by
$h\Psi(x)=k$. Put $\sigma^2=h\mathscr D\Psi(x)$. For fixed real $u$,
the ratio of the centered moment generating functions associated with
$G$ and $F$ is
\[
 \frac{(G/F)(x\mathrm e^{u/\sigma};\mathbf L)}
 {(G/F)(x;\mathbf L)}.
\]
It tends to one, uniformly in $\mathbf L$. To see this, the endpoint
expansions \eqref{eq:TL36} and the quadratic inverse at $x_c$ show that
\[
 \sup_{0<x<x_c}\frac{xY'(x)}{\sqrt{\mathscr D\Psi(x)}}<\infty.
\]
On the tilted interval $x\mathrm e^{u/\sigma}$ the model variance is
comparable to its value at $x$, and the interval remains subcritical.
Therefore $Y(x\mathrm e^{u/\sigma})-Y(x)=O_u(h^{-1/2})$.
Now use \eqref{eq:boundary-real} and the bounded derivative of
$\mathfrak b_Y$.

The real Gaussian limit \eqref{eq:TL39} thus holds for $G$ with the same
center $k$ and scale $\sigma$. Its Fourier bound
\eqref{eq:boundary-fourier} is identical to the one used in
\eqref{eq:TL40}--\eqref{eq:TL46}. Fourier inversion consequently gives
the local limit \eqref{eq:TL47} for $G$ as well. At the common index
$k$, both local probabilities are asymptotic to
$(\sqrt{2\pi}\sigma)^{-1}$. Their ratio tends to one, and hence
\[
 \frac{[x^k]G(x;\mathbf L)}{[x^k]F(x)}
 =\prod_i\mathfrak b_{Y}(L_i)(1+o(1)),\qquad
 \lambda(Y)=\frac{2(n-p)}{2g-2+p}.
\]
All errors are uniform on the fixed length box, including along
sequences tending to either endpoint of the parameter interval.

For bounded $k$, pass to a subsequence on which $k$ is fixed.
Equation~\eqref{eq:TL23} at $Y=0$, with $p+k$ markings, gives the
fixed-cusp boundary estimate \eqref{eq:boundary-factor}. This yields the
same conclusion since the saddle then tends to zero.

Finally, replacing $2(n-p)/(2g-2+p)$ by $n/g$ changes the corresponding
$Y$ by $o(1)$ uniformly in $n\ge p$. Indeed, along a subsequence on
which $n/g$ tends to zero, a positive finite limit, or infinity, both
parameters have the same limit in $[0,y_0]$. Continuity of
$\mathfrak b_Y(L)$, uniformly on the length box, proves
\eqref{eq:boundary-uniform}. At $Y=0$, $T=\Delta=\pi^2/4$;
expanding $I_0$ and integrating gives the first formula in
\eqref{eq:boundary-endpoints}. At $Y=y_0$, the integrand is constant,
which gives the second.
\end{proof}

